% ============================================================
% 08_limits.tex
% Limitations, Ethics Statement, Conclusion
% ============================================================

\section{Limitations}
\label{sec:limitations}

Five limitations bound our claims. First, the selection confound. Users choose their platforms, and the same-event design controls topic and time, not who is posting. Any non-invariance we detect therefore bundles population composition with measurement mode; because the two can offset as well as compound, the bundle cannot be read as a bound on the pure mode effect, only as the platform-level difference in how the construct is measured \citep{sen2021}. The convergence criterion \citep{czarnek2022} removes classifier artifacts from this bundle; it cannot remove selection. Separating the two would require within-user cross-posting data, which we do not have.

Second, platform scope. We study four platforms, and the anchor-density audit may narrow the analysis to the subset with sufficient same-event coverage. All findings are conditional on the retained platforms and events. Exclusion of a sparse platform is a data limitation, not evidence that it measures equivalently. Likewise, following \citet{robitzsch2023}, we read invariance results as diagnostic evidence about this corpus, not as a general licence to pool.

Third, construct. Our instrument measures expressed sentiment, not policy support. Negative sentiment in a thread about inflation need not be opposition to any policy \citep{oconnor2010}. Our equivalence results concern the sentiment measure itself; any mapping to attitudes would need separate validation.

Fourth, a single country and language. All data are UK-focused and English. Machine scoring is language-dependent \citep{nogara2025}, and platform affordance mixes differ across national ecologies, so we make no cross-country claim.

Finally, provenance. The dissertation underlying this paper reported several internally inconsistent statistics: per-platform counts summing to roughly 395,000 against a 279,000 headline (pp.~21-25), a forecast MAE of approximately 0.02 in the text but 0.05 in the goal table (pp.~39, 46-47), post-calibration ECE of 0.05 in the methodology but 0.03 in the results (pp.~30, 36-37), manual labelling coverage of 0.07\% in one chapter and approximately 0.5\% in another (pp.~25, 45), an accuracy figure appearing as both 88.06\% and approximately 0.89 (pp.~35, 50), and a forecast figure titled with a platform the study excluded (p.~38). \todo{List corrected totals and recomputed values.} All statistics here were recomputed from the raw archived data, and the recomputed figures supersede the dissertation throughout. \todo{State the retention basis for the raw archive: the dissertation committed to deleting raw data after project completion and retaining only anonymised aggregates (pp.~15, 25), which the recomputation-from-raw-data claim must reconcile.}

\section{Ethics Statement}
\label{sec:ethics}

This is an observational study of public posts on public policy topics. Collection was approved under the University of Portsmouth ethics review process (School of Computing, Faculty Ethics Committee route; certificate TETHIC-2025-111094, application 27 May 2025, supervisor sign-off 1 June 2025), which classified the project as low risk and complies with GDPR: no private, deleted, or direct-message content was collected, released artifacts contain no personal identifiers, and quoted examples are paraphrased. The dissertation-stage criterion set comprised 200 manually labelled posts balanced across platforms, with inter-annotator agreement of approximately 0.78 (Cohen's kappa; dissertation p.~28). Annotators were \todo{recruitment and compensation details}. The main foreseeable misuse is reading our diagnostics as a ranking of platforms by accuracy; they are equivalence tests, not quality scores, and we caution against that reading.

\section{Conclusion}
\label{sec:conclusion}

Platform is part of the measurement instrument until shown otherwise. We formalise that claim and test it with the psychometric equivalence toolkit, applied to machine-labelled policy sentiment with platform as the grouping facet, to our knowledge for the first time. Either outcome is informative: invariance would give pooled cross-platform indices a defensible footing; non-invariance would quantify exactly where, and by how much, pooling misleads. \todo{One-sentence headline finding once Step 3 and Step 6 results are in.}
